\documentclass[10pt,a4paper]{amsart}
\usepackage[margin=19mm]{geometry}
\usepackage{amsmath,amssymb,mathtools}
\usepackage[scr=rsfs,cal=euler]{mathalfa}
\usepackage{xfrac}
\usepackage[unicode,hidelinks,bookmarks=false]{hyperref}
\newcommand{\ie}{i.e.\ }
\newcommand{\cf}{cf.\ }
\newcommand{\resp}{respectively\ }
\newcommand{\Subsection}{\subsection}
\providecommand{\nobreakdash}{\nobreak\hbox{-}}
\setlength{\emergencystretch}{2em}
\multlinegap0pt
\usepackage{amssymb}
\usepackage{mathtools}
\usepackage[shortlabels]{enumitem}
\newtheorem{lemma}{Lemma}
\newtheorem{theorem}{Theorem}
\title{Characterization of continuous stationary fields as generalized Ornstein-Uhlenbeck fields via multi-parameter Langevin equation and multiple Riemann-Stieltjes integration}
\author{Marko Voutilainen, Pauliina Ilmonen, Lauri Viitasaari}
\date{Source: arXiv:2511.11110v2 (CC BY 4.0)}
\begin{document}
\maketitle

\noindent\emph{Source and licence.} Characterization of continuous stationary fields as generalized Ornstein-Uhlenbeck fields via multi-parameter Langevin equation and multiple Riemann-Stieltjes integration. Version arXiv:2511.11110v2 (CC BY 4.0), distributed by arXiv under the Creative Commons Attribution 4.0 International licence, http://creativecommons.org/licenses/by/4.0/, as recorded in the arXiv licence metadata for this exact version and shown on the abstract page https://arxiv.org/abs/2511.11110v2. Full licence notice: \texttt{licenses/arxiv-2511.11110-CC-BY-4.0.txt}.\par\smallskip

\noindent\textit{Editorial scope.} Counted unit: the article's lemma lemma:ipp (source0.tex lines 355--371). The article's complete original proof of it is delivered with it (source0.tex lines 572--643, one \texttt{proof} environment of 72 lines, which the article places away from the statement). No mathematics is written, repaired or re-derived: the statement and the proof are the article's own lines, in the article's own notation, and no retained line is substituted or re-flowed.  Retained uncounted as the source-local context the proof needs: lines 144--150; lines 262--265; lines 328--334; lines 473--499; lines 567--570.  Nothing was omitted from any retained span, so the package carries no omission record.  No retained line contains a citation, so the wrapper carries no bibliography and no bibliography entry.  No retained line contains a figure environment, an image-inclusion command or drawing code, so no image is delivered and the package declares no asset dependency.  Editorial formatting only, in addition: an AMS \texttt{amsart} preamble that restates the article's own declarations and macros used by the retained lines, namely the environments \texttt{lemma}, \texttt{theorem} on separate counters, together with the standard packages those lines need.  The retained environments print in this document as Lemma 1, Theorem 1 in order of appearance, whereas the article prints the same items with the numbering of its own full document.  The complete native source of the article as distributed in the arXiv e-print for this exact version is bundled unchanged as \texttt{sources/189095/source0.tex}.
\begin{lemma}
\label{binomial}
Let $M\in\mathbb{N}^+$. Then
\begin{equation*}
\sum_{m=0}^M (-1)^m\binom{M}{m} = 0.
\end{equation*}
\end{lemma}

\begin{equation}
\label{rssum}
\lim_{\Vert P \Vert\to 0} \sum_{i_1 = 1}^{n_1}\dots\sum_{i_N = 1}^{n_N} g(\xi_i)[f]_{(x^{(1)}_{i_1-1}, \dots, x^{(N)}_{i_N-1})}^{(x^{(1)}_{i_1}, \dots, x^{(N)}_{i_N})}
\end{equation}

\begin{lemma}
\label{lemma:variationofdifferentiable}
Let $f$ be a function such that $f_t$ exists and is continuous on $[s,t]$. Then
\begin{equation*}
V_{HK}(f, [s,t]) = \sum_{\emptyset\neq v\subseteq \{1,\dots,N\}} \int_{s_v}^{t_v} \left|f_{t_v}(r_v:t_{-v})\right| dr_v.
\end{equation*}
\end{lemma}

\begin{theorem}
\label{lemma:radon}
Let $f$ be a function such that $f_t$ exists and is continuous on $[s,t]$. Then
$$d_u f(r) = f_{t_u}(r)dr_u\quad \text{for all } u\subseteq\{1,\dots,N\}$$
in the multiple Riemann-Stieltjes sense for every continuous integrand $g$. 
\begin{proof}
It suffices to consider the case $u = \{1,\dots,N\}$. The result for lower dimensional differentials follow by fixing variables $r_{-u}$ and considering $f$ as a $|u|$-variate function. Since the integrand $g$ is continuous, both integrals, Riemann-Stieltjes on the left and Riemann on the right, exist. Denote by
$$\left|I_{i_1,\dots,i_N}\right| = \prod_{j=1}^N (x_{i_j}^{(N)} -x_{i_j-1}^{(N)}) $$
the Lebesgue measure of $I_{i_1,\dots,i_N}$. We examine a summand of the Riemann-Stieltjes sum \eqref{rssum}. By the mean value theorem
\begin{equation*}
\begin{split}
 g(\xi_i)[f]_{(x^{(1)}_{i_1-1}, \dots, x^{(N)}_{i_N-1})}^{(x^{(1)}_{i_1}, \dots, x^{(N)}_{i_N})} =
 g(\xi_i) f_{t}(\eta_{i})\left|I_{i_1,\dots,i_N}\right|
 \end{split}
\end{equation*}
with some $\eta_{i}\in I_{i_1,\dots,i_N}$. 
Since $g$ is Riemann-Stieltjes integrable with respect to $f$, we may select $\xi_i = \eta_{i}$
yielding
\begin{equation*}
\begin{split}
&\lim_{\Vert P \Vert\to 0} \sum_{i_1 = 1}^{n_1}\dots\sum_{i_N = 1}^{n_N}g(\xi_i)[f]_{(x^{(1)}_{i_1-1}, \dots, x^{(N)}_{i_N-1})}^{(x^{(1)}_{i_1}, \dots, x^{(N)}_{i_N})}\\
&= \lim_{\Vert P \Vert\to 0}\sum_{i_1 = 1}^{n_1}\dots\sum_{i_N = 1}^{n_N} g(\eta_i) f_{t}(\eta_i)\left|I_{i_1,\dots,i_N}\right|,
\end{split}
\end{equation*}
where the right-hand side converges to the Riemann integral of $gf_t$.
\end{proof}
\end{theorem}

\begin{equation}
\label{productrule}
d(f(r)g(r)) = \sum_{u\subseteq \{1,\dots,N\}} f_{t_u}(r) d_{-u}g(r) dr_u
\end{equation}

\begin{lemma}
\label{lemma:ipp}
Let $f$ and $g$ be functions defined on $[s,t]$. If
\begin{equation*}
\int_{s_v}^{t_v}[g(r) d_vf(r)]_{s_{-v}}^{t_{-v}}
\end{equation*}
exists for all $\emptyset \neq v \subseteq \{1,\dots,N\}$, then also
\begin{equation*}
\int_{s_v}^{t_v}[f(r) d_vg(r)]_{s_{-v}}^{t_{-v}}
\end{equation*}
exists for all $\emptyset \neq v \subseteq \{1,\dots,N\}$. In addition, the integration by parts formula 
\begin{equation}
\label{iip}
\int_s^t f(r) dg(r) = [f(r)g(r)]_s^t + \sum_{\emptyset\neq v \subseteq \{1,\dots,N\}}(-1)^{|v|} \int_{s_v}^{t_v} [g(r)d_vf(r)]_{s_{-v}}^{t_{-v}}
\end{equation}
holds.
\end{lemma}

\begin{proof}
By integration by parts and Theorem \ref{lemma:radon}, the left-hand side yields
\begin{equation*}
\begin{split}
&\int_s^t h(r) d(f(r)g(r)) = [h(r)f(r)g(r)]_s^t + \sum_{\emptyset\neq v \subseteq \{1,\dots,N\}}(-1)^{|v|} \int_{s_v}^{t_v} [g(r)f(r)d_vh(r)]_{s_{-v}}^{t_{-v}}\\
&= [h(r)f(r)g(r)]_s^t +\sum_{\emptyset\neq v\subseteq \{1,\dots,N\}}(-1)^{|v|} \int_{s_v}^{t_v} [g(r)f(r)h_{t_v}(r)dr_v]_{s_{-v}}^{t_{-v}}.
 \end{split}
\end{equation*}
We apply Theorem \ref{lemma:radon} and Lemma \ref{lemma:ipp} repeatedly in the rest of the proof.
The right-hand side of \eqref{productrule} gives
\begin{equation*}
\begin{split}
\int_s^t h(r)f(r) dg(r) +  \sum_{\emptyset\neq u\subseteq \{1,\dots,N\}}\int_s^t  h(r)f_{t_u}(r) d_{-u}g(r) dr_u,
\end{split}
\end{equation*}
where
\begin{equation*}
\int_s^t h(r)f(r) dg(r) = [h(r)f(r)g(r)]_s^t + \sum_{\emptyset\neq v \subseteq \{1,\dots,N\}}(-1)^{|v|} \int_{s_v}^{t_v} [g(r)d_v(h(r)f(r))]_{s_{-v}}^{t_{-v}}
\end{equation*}
and
\begin{equation*}
\begin{split}
&\sum_{\emptyset\neq u\subseteq \{1,\dots,N\}}\int_s^t  h(r)f_{t_u}(r) d_{-u}g(r) dr_u = \sum_{\emptyset\neq u\subseteq \{1,\dots,N\}}\int_{s_u}^{t_u} \int_{s_{-u}}^{t_{-u}} h(r)f_{t_u}(r) d_{-u}g(r) dr_u \\
&=\int_{s}^{t} h(r)f_t(r)g(r) dr + \sum_{\substack{u\subseteq\{1,\dots,N\}\\
1 \leq |u|\leq N-1}} \int_{s_u}^{t_u} [h(r)f_{t_u}(r)g(r)]_{s_{-u}}^{t_{-u}} dr_u\\
&\quad+ \sum_{\substack{u\subseteq\{1,\dots,N\}\\
1 \leq |u|\leq N-1}}\int_{s_u}^{t_u} \sum_{\substack{v\subseteq -u\\ 1\leq |v|}} (-1)^{|v|} \int_{s_v}^{t_v} [g(r) d_v(h(r)f_{t_u}(r))]_{s_{-u-v}}^{t_{-u-v}} dr_u.
\end{split}
\end{equation*}
Note that the integrals in the last term are well-defined by Lemma \ref{lemma:variationofdifferentiable}, since the partial derivative of $hf_{t_u}$ with respect to $t_v$ is continuous. By combining the previous two equations, the right-hand side of \eqref{productrule} gives
\begin{equation*}
\begin{split}
&[h(r)f(r)g(r)]_s^t +\sum_{\emptyset\neq v \subseteq \{1,\dots,N\}}(-1)^{|v|} \int_{s_v}^{t_v} [g(r)\frac{dhf}{dt_v}(r) dr_v]_{s_{-v}}^{t_{-v}}\\
&\quad+\int_{s}^{t} h(r)f_t(r)g(r) dr + \sum_{\substack{u\subseteq\{1,\dots,N\}\\
1 \leq |u|\leq N-1}} \int_{s_u}^{t_u} [h(r)f_{t_u}(r)g(r)dr_u]_{s_{-u}}^{t_{-u}}\\
&\quad+ \sum_{\substack{u\subseteq\{1,\dots,N\}\\
1 \leq |u|\leq N-1}} \sum_{\substack{v\subseteq -u\\ 1\leq |v|}} (-1)^{|v|} \int_{s_{u+v}}^{t_{u+v}} \left[g(r) \frac{dhf_{t_u}}{dt_v}(r) dr_{u+v}\right]_{s_{-u-v}}^{t_{-u-v}}.
\end{split}
\end{equation*}
That is, in order to show that Equation \eqref{productrule} holds in the given sense, we need to show that
\begin{equation*}
\begin{split}
&\sum_{\emptyset\neq v \subseteq \{1,\dots,N\}}(-1)^{|v|} \int_{s_v}^{t_v} [g(r)\sum_{\substack{u\subseteq v\\ u \neq v}}h_{t_u}(r)f_{t_{v-u}}(r) dr_v]_{s_{-v}}^{t_{-v}}\\
&\quad +\int_{s}^{t} h(r)f_t(r)g(r) dr + \sum_{\substack{u\subseteq\{1,\dots,N\}\\
1 \leq |u|\leq N-1}} \int_{s_u}^{t_u} [h(r)f_{t_u}(r)g(r)dr_u]_{s_{-u}}^{t_{-u}}\\
&\quad+ \sum_{\substack{u\subseteq\{1,\dots,N\}\\
1 \leq |u|\leq N-1}} \sum_{\substack{v\subseteq -u\\ 1\leq |v|}} (-1)^{|v|} \int_{s_{u+v}}^{t_{u+v}} \left[g(r) \sum_{w\subseteq v} h_{t_w}(r) f_{t_{u+(v-w)}}(r) dr_{u+v}\right]_{s_{-u-v}}^{t_{-u-v}}\\
&= \sum_{\emptyset\neq v \subseteq \{1,\dots,N\}}(1+(-1)^{|v|}) \int_{s_v}^{t_v} [g(r)h(r)f_v(r) dr_v]_{s_{-v}}^{t_{-v}}\\
&\quad+ \sum_{\substack{u\subseteq\{1,\dots,N\}\\
1 \leq |u|\leq N-1}} \sum_{\substack{v\subseteq -u\\ 1\leq |v|}} (-1)^{|v|} \int_{s_{u+v}}^{t_{u+v}} \left[g(r) h(r) f_{t_{u+v}}(r) dr_{u+v}\right]_{s_{-u-v}}^{t_{-u-v}}\\
&\quad+ \sum_{\substack{v\subseteq \{1,\dots,N\} \\ 2 \leq |v|}}(-1)^{|v|} \int_{s_v}^{t_v} [g(r)\sum_{\substack{u\subseteq v\\ 1\leq |u| \leq |v|-1}}h_{t_u}(r)f_{t_{v-u}}(r) dr_v]_{s_{-v}}^{t_{-v}}\\
&\quad+ \sum_{\substack{u\subseteq\{1,\dots,N\}\\
1 \leq |u|\leq N-1}} \sum_{\substack{v\subseteq -u\\ 1\leq |v|}} (-1)^{|v|} \int_{s_{u+v}}^{t_{u+v}} \left[g(r) \sum_{\substack{w\subseteq v\\1\leq |w|}} h_{t_w}(r) f_{t_{u+(v-w)}}(r) dr_{u+v}\right]_{s_{-u-v}}^{t_{-u-v}} =0.
\end{split}
\end{equation*}
The first term after the equality follows by choosing $u = \emptyset$ in the first term of the equation and combining it with the following two terms. The third term after the equality is the first term of the equation with the additional restriction $u \neq \emptyset$. 
We show that the sum of the last two terms equals zero. The first two terms can be handled in a similar, but less involved manner. Thereby, it suffices to show
\begin{equation}
\label{zero}
\begin{split}
&\sum_{\substack{q\subseteq \{1,\dots,N\} \\ 2 \leq |q|}}(-1)^{|q|} \int_{s_q}^{t_q} [g(r)\sum_{\substack{w\subseteq q\\ 1\leq |w| \leq |q|-1}}h_{t_w}(r)f_{t_{q-w}}(r) dr_q]_{s_{-q}}^{t_{-q}}\\
&\quad+ \sum_{\substack{u\subseteq\{1,\dots,N\}\\
1 \leq |u|\leq N-1}} \sum_{\substack{v\subseteq -u\\ 1\leq |v|}} (-1)^{|v|} \int_{s_{u+v}}^{t_{u+v}} \left[g(r) \sum_{\substack{\tilde{w}\subseteq v\\1\leq |\tilde{w}|}} h_{t_{\tilde{w}}}(r) f_{t_{u+(v-\tilde{w})}}(r) dr_{u+v}\right]_{s_{-u-v}}^{t_{-u-v}} = 0.
\end{split}
\end{equation}
Now, let $q \subseteq \{1,\dots,N\}$ with $|q|\geq 2$ be fixed. Also, let $w\subseteq q$ with $1\leq |w|\leq |q|-1$ be fixed. By doing this, we have fixed a unique summand of the first term of \eqref{zero}. The sign of this term is positive if $|q|$ is even and negative if $|q|$ is odd. We next consider which terms in the latter double sum of \eqref{zero} correspond to the choices $q$ and $w$. First of all, it has to hold that $w\subseteq v$, but $v$ can also contain indices of $q-w$. In addition, $u+v=q$ has to hold. That is, by fixing $v$ we also fix $u$. We now have four different scenarios, where $|q|$ and $|w|$ can be even or odd. We start by assuming that $|q|$ and $|w|$ are both even. By Lemma \ref{binomial}, the sum related to the corresponding terms in \eqref{zero} equals
\begin{equation}
\label{zero1}
\binom{|q|-|w|}{0} - \binom{|q|-|w|}{1}+\dots -\binom{|q|-|w|}{|q|-|w|-1} + 1 = 0,
\end{equation}
where the binomial coefficient $\binom{|q|-|w|}{i}$ corresponds to $v$ which contains $w$ and $i$ elements from $q-w$. Here we have used the fact that since $u+v = q$ and $|u|,|v| \geq 1$, $i$ runs from 0 to $|q| - |w| -1$. By considering in a similar way the other cases, we see that regardless of whether $|q|$ and $|w|$ are even or odd, the sum equals zero by Lemma \ref{binomial}. 
\end{proof}

\end{document}
